\uuid{ARrV}
\titre{Espace des fonctions}
\auteur{hanani ; mzari}
\organisation{exo7}
\licence{CC BY-NC-SA 4.0}
\datecreate{2020-01-20}
\niveau{L1}
\module{Algèbre}
\chapitre{Espace vectoriel}
\sousChapitre{Définition, sous-espace}
\difficulte{3}
\theme{Espaces vectoriels, Espace des fonctions}
\exo7qcmid{lille-2.espaces-vectoriels.016}

\contenu{
\begin{question}
\qclassification{106.01}
\qtags{motcle=Espace des fonctions}
Soit $E=\{f:\Rr \to \Rr \, ; \; f\mbox{ est croissante sur }\Rr\}$.
Quelles sont les assertions vraies ?
\begin{answers}
\good{La fonction nulle appartient à $E$.}
\good{$E$ est stable par addition.}
\bad{$E$ est stable par multiplication par un scalaire.}
\bad{$E$ est un espace vectoriel.}
\end{answers}
\begin{explanations} La fonction nulle appartient à $E$, puisqu'elle est constante.\\
On vérifie que $E$ est stable par addition. Par contre, $E$ ne l'est pas par multiplication par un scalaire $<0$. Donc
$E$ n'est pas un espace vectoriel.
\end{explanations}
\end{question}
}
